\documentclass[11pt,a4paper,fontset=fandol]{ctexart}

\usepackage[a4paper,top=2.25cm,bottom=2.45cm,left=2.25cm,right=2.25cm,headheight=18pt,headsep=0.55cm,footskip=24pt]{geometry}
\usepackage{amsmath,amssymb,amsthm}
\usepackage{fancyhdr}
\usepackage{lastpage}
\usepackage{enumitem}
\usepackage{titlesec}
\usepackage{xcolor}
\usepackage{hyperref}
\usepackage{bookmark}
\usepackage[most]{tcolorbox}
\usepackage{setspace}
\usepackage{array}

\hypersetup{
  colorlinks=true,
  linkcolor=blue!50!black,
  urlcolor=blue!50!black
}

\setstretch{1.13}
\setlength{\parindent}{2em}
\setlength{\parskip}{0.25em}
\allowdisplaybreaks
\setlist[enumerate]{itemsep=0.45em, topsep=0.35em, leftmargin=2.4em}
\setlist[itemize]{itemsep=0.25em, topsep=0.25em, leftmargin=1.9em}

\definecolor{mainblue}{RGB}{36,83,140}
\definecolor{lightblue}{RGB}{241,246,252}
\definecolor{lightgray}{RGB}{246,246,246}
\definecolor{rulegray}{RGB}{110,118,128}

\titleformat{\section}
  {\Large\bfseries\color{mainblue}}
  {\thesection.}
  {0.45em}
  {}
  [\vspace{0.2em}\color{rulegray}\titlerule]
\titlespacing*{\section}{0pt}{1.1em}{0.7em}

\pagestyle{fancy}
\fancyhf{}
\fancyhead[L]{\small 图论计数周测 B 卷}
\fancyhead[C]{\small 组合专题：结构证明与路径变式}
\fancyhead[R]{\small 刘文博}
\fancyfoot[C]{\small 第 \thepage 页 / 共 \pageref{LastPage} 页}
\renewcommand{\headrulewidth}{0.55pt}
\renewcommand{\footrulewidth}{0.35pt}

\newtcolorbox{infobox}[1]{
  breakable,
  colback=lightblue,
  colframe=mainblue,
  boxrule=0.7pt,
  arc=1mm,
  title=\textbf{#1},
  fonttitle=\bfseries
}

\newenvironment{solution}{\par\noindent\textbf{参考解答：}}{\hfill$\square$\par}
\newcommand{\answerspace}{\vspace{2.3cm}}
\newcommand{\biganswerspace}{\vspace{3.2cm}}

\begin{document}

\thispagestyle{empty}
\begin{center}
{\fontsize{22}{28}\selectfont\bfseries 图论计数周测 B 卷\par}
\vspace{0.4cm}
{\large 适合小学高年级至初中低年级数学竞赛训练\par}
\end{center}

\begin{infobox}{考试说明}
\begin{itemize}
  \item 测试时间：75 分钟
  \item 满分：100 分
  \item B 卷更强调建模、结构证明和网格路径变式，建议先判断是用握手定理、完全图/二分图，还是用路径分段计数。
\end{itemize}
\end{infobox}

\section{试题}

\begin{enumerate}
  \item （12 分）一个图有 12 个点，每个点的度数都等于 5。求这个图的边数。
  \answerspace

  \item （12 分）11 个同学参加循环赛，每两人比赛 1 场，共有多少场比赛？
  \answerspace

  \item （12 分）从 $(0,0)$ 到 $(5,3)$，每次只能向右或向上走 1 格，共有多少条最短路径？
  \answerspace

  \item （12 分）从 $(0,0)$ 到 $(4,3)$，每次只能向右或向上走 1 格，但不允许经过点 $(2,2)$。共有多少条最短路径？
  \biganswerspace

  \item （12 分）证明：任意一个图中，奇度数点的个数一定是偶数。
  \answerspace

  \item （12 分）一个图有 8 个点，且每个点的度数都不小于 4。证明这个图的边数至少为 16。
  \answerspace

  \item （14 分）判断下列命题真假，并说明理由：在任意一个图中，度数至少为 2 的点不可能恰好只有 1 个。
  \biganswerspace

  \item （14 分）在任意 6 个人中，一定可以找到 3 个人，他们之间两两认识，或者 3 个人，他们之间两两不认识。请用图论语言表述并证明这个结论。
  \biganswerspace
\end{enumerate}

\newpage
\section{详细解答}

\begin{enumerate}
  \item \begin{solution}
  总度数为
  \[
  12\times 5=60.
  \]

  由握手定理
  \[
  2E=60,
  \]
  所以
  \[
  E=30.
  \]
  \end{solution}

  \item \begin{solution}
  11 个同学两两比赛 1 场，就是完全图 $K_{11}$ 的边数，因此总场数为
  \[
  \binom{11}{2}=55.
  \]
  \end{solution}

  \item \begin{solution}
  从 $(0,0)$ 到 $(5,3)$，需要走 5 步向右和 3 步向上，共 8 步。

  所以最短路径条数为
  \[
  \binom83=56.
  \]
  \end{solution}

  \item \begin{solution}
  先求总最短路径数：
  \[
  \binom73=35.
  \]

  再求经过 $(2,2)$ 的最短路径数。

  从 $(0,0)$ 到 $(2,2)$，需走 2 步向右、2 步向上，所以有
  \[
  \binom42=6
  \]
  条。

  从 $(2,2)$ 到 $(4,3)$，需走 2 步向右、1 步向上，所以有
  \[
  \binom31=3
  \]
  条。

  因此经过 $(2,2)$ 的最短路径共有
  \[
  6\times 3=18
  \]
  条。

  所以不经过 $(2,2)$ 的最短路径数为
  \[
  35-18=17.
  \]
  \end{solution}

  \item \begin{solution}
  由握手定理，图中所有点的度数之和等于边数的 2 倍，因此总度数和一定是偶数。

  若一个和是偶数，那么其中奇数加数的个数必须是偶数；否则奇数个奇数相加会得到奇数。

  所以任意一个图中，奇度数点的个数一定是偶数。
  \end{solution}

  \item \begin{solution}
  因为每个点的度数都不小于 4，所以总度数至少为
  \[
  8\times 4=32.
  \]

  由握手定理
  \[
  2E\ge 32,
  \]
  所以
  \[
  E\ge 16.
  \]
  \end{solution}

  \item \begin{solution}
  这个命题是\textbf{假的}。

  反例：取 3 个点，画成一个“V”字形，也就是中间点分别与另外两个点相连，而这两个端点之间不连边。

  此时三个点的度数分别为
  \[
  2,\ 1,\ 1.
  \]

  于是度数至少为 2 的点恰好只有 1 个。

  所以命题不成立。
  \end{solution}

  \item \begin{solution}
  把 6 个人看成 6 个点；如果两个人认识，就在对应两点之间连红边；如果两个人不认识，就连蓝边。

  这样任意两点之间恰有一条红边或蓝边。

  任取其中一个点 $A$，其余 5 个点与 $A$ 的连边只有两种颜色：红或蓝。由抽屉原理，这 5 条边中至少有 3 条颜色相同。

  不妨设 $A$ 与 $B,C,D$ 三点之间都是红边。

  \begin{itemize}
    \item 如果 $B,C,D$ 中有两点之间也是红边，例如 $B$ 与 $C$ 之间是红边，那么 $A,B,C$ 就是 3 个两两认识的人；
    \item 如果 $B,C,D$ 中任意两点之间都不是红边，那么它们之间就都必须是蓝边，于是 $B,C,D$ 是 3 个两两不认识的人。
  \end{itemize}

  两种情况总有一种发生，因此结论成立。
  \end{solution}
\end{enumerate}

\end{document}
